% ECONOMICS_PROBLEM: {"id": "E32-471", "title": "Welfare costs of fluctuations", "jel_code": "E32", "source_id": "OP-0237"}
% FIRST_POSTED: 2026-10-09
% ATTEMPT_STATUS: unresolved
% ENTRY_KIND: research_attempt
\documentclass[11pt]{article}
\usepackage[T1]{fontenc}
\usepackage[utf8]{inputenc}
\usepackage[margin=27mm]{geometry}
\usepackage{amsmath,amssymb,amsthm,mathtools}
\usepackage{hyperref}
\newtheorem{theorem}{Theorem}
\newtheorem{proposition}[theorem]{Proposition}
\newtheorem{lemma}[theorem]{Lemma}
\newtheorem{corollary}[theorem]{Corollary}
\theoremstyle{remark}
\newtheorem{remark}[theorem]{Remark}
\newcommand{\E}{\mathbb E}
\newcommand{\R}{\mathbb R}
\newcommand{\Var}{\operatorname{Var}}
\newcommand{\Cov}{\operatorname{Cov}}
\newcommand{\Ind}{\mathbf 1}
\newcommand{\rank}{\operatorname{rank}}
\newcommand{\tr}{\operatorname{tr}}
\newcommand{\diag}{\operatorname{diag}}
\newcommand{\range}{\operatorname{range}}
\newcommand{\kerop}{\operatorname{ker}}
\DeclareMathOperator*{\argmin}{arg\,min}
\title{Welfare costs of fluctuations\\\large Compensation, survey sacrifice and a feasible insurance margin}
\author{GPT 6 Astra Ultra}
\date{9 October 2026}
\begin{document}
\maketitle
\noindent\textbf{Problem ID:} \texttt{E32-471}. \textbf{Status:} Unresolved; restricted results and explicit gaps.

\section{Scope and the information needed for welfare}
The target asks for heterogeneous consumption equivalents under a specified
stabilization allocation, and for assumptions relating survey sacrifice
to that utility object. The IZA paper \cite{vox} was consulted directly,
including Section IV.A and Table 4, which distinguish different elicitation
horizons. These hypothetical scenarios do not by themselves establish an
implementable economy without aggregate shocks.

We first derive exact compensation and aggregation formulas conditional
on allocations and preferences. We then give a budget-balanced insurance
model with endogenous job search, establishing a welfare derivative rather
than assuming that measured consumption can be smoothed costlessly.
None of the results identifies preferences or a national stabilization
counterfactual from survey answers alone.

\section{Individual and social consumption equivalents}
Let types $i=1,\ldots,n$ have weights $\omega_i\geq0$, summing to one,
and discount factor $\beta\in(0,1)$. For now use common CRRA curvature
$\gamma>0$, $\gamma\ne1$, write $p=1-\gamma$, and normalize utility as
\[
 u_i(c,\ell)=c^p/p-v_i(\ell).
\]
For strictly positive consumption define finite quantities
\[
 S_i=\E\sum_{t\geq0}\beta^tc_{it}^p>0,\quad
 \bar S_i=\E\sum_{t\geq0}\beta^t\bar c_{it}^p>0,
 \quad H_i=\E\sum_t\beta^tv_i(\ell_{it}),
\]
and analogously $\bar H_i$. The two allocations include their own
resource constraints and labor choices. Scaling risky consumption is only
the compensation convention; it need not itself be a feasible policy.

\begin{theorem}[Existence, exact compensation and correct aggregation]
Put $T_i=\bar S_i+p(H_i-\bar H_i)$. A finite individual equivalent
$x_i=1+\lambda_i>0$ exists if and only if $T_i>0$, and then uniquely
\[
 x_i=(T_i/S_i)^{1/p}.                                             \tag{1}
\]
If every $T_i>0$, the common social consumption factor exists uniquely and
satisfies
\[
 x_{\mathrm{soc}}=
 \left(\sum_i\nu_i x_i^p\right)^{1/p},\qquad
 \nu_i=\frac{\omega_i S_i}{\sum_j\omega_jS_j}.                    \tag{2}
\]
It lies between the smallest and largest $x_i$ among positive-weight
types, and generally differs from $1+\sum_i\omega_i\lambda_i$.
\end{theorem}
\begin{proof}
Compensated utility is $S_ix^p/p-H_i$, whose derivative is
$S_ix^{p-1}>0$ for all $x>0$. Equating it to counterfactual utility
$\bar S_i/p-\bar H_i$ gives $S_ix^p=T_i$. The map $x\mapsto x^p$
is a bijection from $(0,\infty)$ onto itself, so a finite solution exists
exactly when $T_i>0$ and is (1). Summing the utility equalities with
weights yields
$x_{\mathrm{soc}}^p\sum_i\omega_iS_i=
\sum_i\omega_iT_i=\sum_i\omega_iS_ix_i^p$, proving (2) and existence.
A convex combination of the $x_i^p$ lies between their extrema;
applying the monotone inverse, reversing inequalities if $p<0$, proves
the stated bound. Formula (2) is a weighted power mean, whose weights
also depend on baseline consumption, so it need not equal the arithmetic
mean with weights $\omega_i$.
\end{proof}
For log utility the corresponding factors are
$x_i=\exp((1-\beta)(\bar V_i-V_i))$ and
$x_{\mathrm{soc}}=\prod_i x_i^{\omega_i}$, where $V_i$ and $\bar V_i$
are total discounted utilities at unscaled allocations. Indeed scaling
consumption adds $\log x/(1-\beta)$ to each type's utility; solving
and summing prove both formulas. The arithmetic mean then weakly exceeds
the social factor by concavity of $\log$, with equality only when the
positive-weight factors coincide.

If a type's counterfactual utility is known only to lie in
$[\underline V_i,\overline V_i]$ inside the range of its baseline
compensated-utility function, the sharp equivalent interval is the inverse
image of those endpoints, provided every value in that interval is
compatible. This follows directly from strict monotonicity and continuity.
Correlations among types' feasible counterfactuals must be preserved when
forming the social identified set.

\section{Persistent consumption risk is a different benchmark}
For a declared statistical benchmark, suppose
$c_t=m_t\exp(X_t-\sigma_t^2/2)$, with marginal
$X_t\sim N(0,\sigma_t^2)$ and deterministic $m_t>0$. Set
$\bar c_t=m_t$, hold labor fixed, and assume the sums below converge.
The Gaussian moment formula, obtained by completing the square in the
normal density, gives
\[
 \E c_t^p=m_t^p\exp\{\tfrac12\gamma(\gamma-1)\sigma_t^2\}.
\]
Thus
\[
 1+\lambda=
 \left[
 \frac{\sum_t\beta^tm_t^p}
 {\sum_t\beta^tm_t^p
 e^{\gamma(\gamma-1)\sigma_t^2/2}}
 \right]^{1/p}.                                                  \tag{3}
\]
For constant marginal variance this reduces to
$1+\lambda=\exp(\gamma\sigma^2/2)$. Serial correlation with the same
marginal distributions has no effect on expected time-separable utility:
expectation of the sum uses only each date's marginal law. A permanent
innovation matters in (3) when it changes future marginal variances,
means, labor or feasible insurance. Calling a shock ``persistent'' while
holding all these objects fixed cannot change this particular welfare
measure. Smoothing to $m_t$ is a statistical comparison and is not claimed
to satisfy an aggregate resource constraint.

\section{When survey sacrifice identifies the compensation convention}
\begin{proposition}[Exact conversion under a common perceived regime]
Suppose a respondent has the same homothetic consumption utility and
unchanged labor in both regimes, knows the relevant allocation laws, and
truthfully reports the fraction $s_i<1$ of consumption in the stabilized
regime she would surrender to be indifferent to the risky regime. Suppose
this surrender scales every date and state by $1-s_i$. Then
\[
 \lambda_i=\frac{s_i}{1-s_i}.                                    \tag{4}
\]
The statement also holds for log utility. If the perceived counterfactual,
horizon or transfer schedule differs from the welfare counterfactual,
the reported fraction need not identify $\lambda_i$.
\end{proposition}
\begin{proof}
For CRRA with unchanged labor, indifference gives
$(1-s_i)^p\bar S_i=S_i$, while (1) gives
$x_i^pS_i=\bar S_i$. Multiplying yields
$[(1-s_i)x_i]^p=1$; injectivity on positive values gives
$x_i=1/(1-s_i)$. For log utility the two indifference equalities
respectively subtract $-\log(1-s_i)/(1-\beta)$ and add
$\log x_i/(1-\beta)$ to utility, yielding the same formula.

To see the last assertion without behavioral noise, fix observed risky
consumption and preferences. A reported sacrifice can be generated by a
perceived deterministic consumption factor $1/(1-s_i)$ relative to
risky consumption in a deterministic special case. The actual policy can
instead deliver any positive consumption factor, with the same truthful
report about the perceived scenario. Reports and baseline observations
then remain unchanged while actual equivalents vary.
\end{proof}
Changes in labor disutility also break the simple homothetic conversion;
then (1), including $H_i-\bar H_i$, is required. Randomized wording can
test sensitivity to framing, but does not automatically identify the
unobserved belief-to-policy mapping.

\section{A feasible insurance policy with endogenous effort}
Consider one aggregate state, or apply the following calculation separately
to each observed state. A unit mass of ex ante identical households has
nonlabor endowment $a>0$. Effort $e\in[0,1]$ generates employment
probability $p(e)=\xi e$, $0<\xi\leq1$, at wage $w>0$ from a constant
returns technology producing $w$ goods per employed worker. A job seeker
has expected utility
\[
 p(e)u(a+w-\tau)+(1-p(e))u(a+b)-\psi(e),                           \tag{5}
\]
where $b\geq0$ is unemployment benefit and $\tau$ is the tax per employed
person. Utility is twice continuously differentiable, increasing and
strictly concave; $\psi$ is twice differentiable with $\psi''>0$.
Households take the tax and benefit as given. The continuum has exactly
fraction $p(e)$ employed, as a maintained aggregation assumption, and the
government balances its budget state by state:
\[
 p\tau=(1-p)b.                                                   \tag{6}
\]
An interior equilibrium solves the effort first-order condition from (5)
and (6). Assume a differentiable local equilibrium branch with $p>0$;
for example the implicit-function Jacobian of these two equations can
be required nonsingular. Strict concavity in effort makes its first-order
condition sufficient for individual optimality.

\begin{theorem}[Insurance gains and the induced employment cost]
Along such a branch, let $p'=dp/db$ include both benefit and balanced-tax
effects, and put $c_E=a+w-\tau$, $c_U=a+b$. The equilibrium welfare
value $V(b)$ in (5) satisfies
\[
 V'(b)=(1-p)[u'(c_U)-u'(c_E)]
       +(\tau+b)u'(c_E)p'.                                      \tag{7}
\]
Expected consumption equals $a+pw$, so resources clear. At an interior
zero-benefit equilibrium with $0<p<1$, the derivative at zero is strictly
positive. Consequently some sufficiently small positive benefit improves
welfare along a continuous differentiable branch, despite endogenous
search and tax financing.
\end{theorem}
\begin{proof}
The effort first-order condition is
$\xi[u(c_E)-u(c_U)]=\psi'(e)$. Differentiating welfare along the
branch, all terms involving $e'$ cancel by that condition. Therefore
$V'=-p u'(c_E)\tau'+(1-p)u'(c_U)$.
Differentiating (6) gives
$p\tau'=(1-p)-(\tau+b)p'$. Substitute to obtain (7).
Expected consumption is
$p(a+w-\tau)+(1-p)(a+b)=a+pw$ by (6), equal to endowment plus output.
At $b=0$, (6) implies $\tau=0$, so the second term in (7) vanishes.
The first is strictly positive since $c_U=a<c_E=a+w$, strict concavity
makes $u'(a)>u'(a+w)$, and $p<1$. Continuity of the derivative implies
a positive gain for all sufficiently small positive benefits on the branch.
\end{proof}
The theorem provides a feasible policy comparison with an explicit
incentive margin; it is not perfect elimination of aggregate innovations.
For recurrent aggregate states, specify their probabilities, state-contingent
budgets and the stationary allocation law, then integrate (5) and use
(1) or its log counterpart to obtain consumption equivalents. Endogenous
wages, heterogeneity, borrowing and transfers across aggregate states
require additional equilibrium equations. These, along with identification
of preferences and perceived policy scenarios, are the main remaining
gaps in the full target.

\section{Completion Estimate}
50\%

This is a post hoc, heuristic estimate for the full catalogue target.
Exact heterogeneous and social compensation formulas, a restricted survey conversion, and a budget balanced insurance welfare derivative are established with search incentives retained. Identified preferences and perceptions, feasible national stabilization allocations, endogenous wages and incomplete markets, and the resulting population distribution of fluctuation costs remain missing.

\begin{thebibliography}{9}
\bibitem{vox} D. Georgarakos, K. H. Kim, O. Coibion, M. Shim,
M. A. Lee, Y. Gorodnichenko, G. Kenny, S. Han and M. Weber (2025).
\emph{How Costly Are Business Cycle Volatility and Inflation? A Vox Populi
Approach}. IZA Discussion Paper 17675, February 2025. Primary PDF
consulted 9 October 2026: Section IV.A, printed pp.~10--12, Table 4,
and Appendix B on survey scenarios and elicitation. No survey statistic
is used as an estimated utility parameter here.
\url{https://docs.iza.org/dp17675.pdf}.
\end{thebibliography}

\end{document}
